\fancyhead[L]{\small Dritter Beweis}
\SourceLeaf{066}
\begin{center}
\textbf{SATZ ÜBER DIE ZERLEGBARKEIT}\\[8pt]
\textbf{VON POLYNOMEN}\\[8pt]
\textbf{IN REELLE FAKTOREN.}\\[8pt]
DRITTER BEWEIS.\\[8pt]
ERGÄNZUNG ZUR VORHERGEHENDEN ABHANDLUNG.\\[12pt]
\rule{3cm}{0.3pt}
\end{center}
Nachdem die vorhergehende Abhandlung bereits gedruckt war, führten erneute Überlegungen über denselben Gegenstand zu einem neuen Beweis des Satzes. Wie der vorhergehende ist auch dieser rein analytisch, beruht aber auf völlig anderen Grundsätzen und scheint ihm hinsichtlich der Einfachheit bei Weitem vorzuziehen zu sein. Diesem \textit{dritten} Beweis seien daher die folgenden wenigen Seiten gewidmet.

\ThirdArticle{1}
Gegeben sei folgende Funktion der Unbestimmten $x$:
\[
X=x^m+Ax^{m-1}+Bx^{m-2}+Cx^{m-3}+\text{usw.}+Lx+M
\]
in der die Koeffizienten $A,B,C$ usw. bestimmte reelle Größen sind. Es seien $r,\varphi$ weitere Unbestimmte, und wir setzen
\[
\begin{aligned}
&r^m\cos m\varphi+Ar^{m-1}\cos(m-1)\varphi\\
&\qquad+Br^{m-2}\cos(m-2)\varphi\\
&\qquad+Cr^{m-3}\cos(m-3)\varphi+\text{usw.}+Lr\cos\varphi+M=t\\
&r^m\sin m\varphi+Ar^{m-1}\sin(m-1)\varphi\\
&\qquad+Br^{m-2}\sin(m-2)\varphi\\
&\qquad+Cr^{m-3}\sin(m-3)\varphi+\text{usw.}+Lr\sin\varphi=u
\end{aligned}
\]
\SourceLeaf{067}
\[
\begin{aligned}
&mr^m\cos m\varphi+(m-1)Ar^{m-1}\cos(m-1)\varphi\\
&\qquad+(m-2)Br^{m-2}\cos(m-2)\varphi\\
&\qquad+(m-3)Cr^{m-3}\cos(m-3)\varphi+\text{usw.}+Lr\cos\varphi=t'\\
&mr^m\sin m\varphi+(m-1)Ar^{m-1}\sin(m-1)\varphi\\
&\qquad+(m-2)Br^{m-2}\sin(m-2)\varphi\\
&\qquad+(m-3)Cr^{m-3}\sin(m-3)\varphi+\text{usw.}+Lr\sin\varphi=u'\\[6pt]
&mmr^m\cos m\varphi+(m-1)^2Ar^{m-1}\cos(m-1)\varphi\\
&\qquad+(m-2)^2Br^{m-2}\cos(m-2)\varphi\\
&\qquad+(m-3)^2Cr^{m-3}\cos(m-3)\varphi+\text{usw.}+Lr\cos\varphi=t''\\
&mmr^m\sin m\varphi+(m-1)^2Ar^{m-1}\sin(m-1)\varphi\\
&\qquad+(m-2)^2Br^{m-2}\sin(m-2)\varphi\\
&\qquad+(m-3)^2Cr^{m-3}\sin(m-3)\varphi+\text{usw.}+Lr\sin\varphi=u''
\end{aligned}
\]
\[
\frac{(tt+uu)(tt''+uu'')+(tu'-ut')^2-(tt'+uu')^2}{r(tt+uu)^2}=y
\]
Der Faktor $r$ lässt sich offenbar aus dem Nenner der letzten Formel kürzen, da $t',u',t'',u''$ durch ihn teilbar sind. Schließlich sei $R$ eine bestimmte positive Größe, die beliebig ist, jedoch die größte der Größen
\[
mA\BareRoot{2},\quad\BareRoot{(mB\BareRoot{2})},\quad\BareRoot[3]{(mC\BareRoot{2})},\quad\BareRoot[4]{(mD\BareRoot{2})}\quad\text{usw.}
\]
übertrifft, wobei von den Vorzeichen der Größen $A,B,C$ usw. abgesehen wird, also etwaige negative Größen unter ihnen durch positive ersetzt werden. Nach diesen Vorbereitungen behaupte ich, dass $tt'+uu'$ gewiss einen positiven Wert annimmt, wenn man $r=R$ setzt, welchen (reellen) Wert man $\varphi$ auch beilegt.

\textit{Beweis.} Wir setzen
\[
\begin{aligned}
&R^m\cos45^\circ+AR^{m-1}\cos(45^\circ+\varphi)\\
&\qquad+BR^{m-2}\cos(45^\circ+2\varphi)\\
&\quad+CR^{m-3}\cos(45^\circ+3\varphi)+\text{usw.}\\
&\qquad+LR\cos(45^\circ+(m-1)\varphi)+M\cos(45^\circ+m\varphi)=T\\
&R^m\sin45^\circ+AR^{m-1}\sin(45^\circ+\varphi)\\
&\qquad+BR^{m-2}\sin(45^\circ+2\varphi)\\
&\quad+CR^{m-3}\sin(45^\circ+3\varphi)+\text{usw.}\\
&\qquad+LR\sin(45^\circ+(m-1)\varphi)+M\sin(45^\circ+m\varphi)=U\\
&mR^m\cos45^\circ+(m-1)AR^{m-1}\cos(45^\circ+\varphi)\\
&\qquad+(m-2)BR^{m-2}\cos(45^\circ+2\varphi)\\
&\quad+(m-3)CR^{m-3}\cos(45^\circ+3\varphi)+\text{usw.}\\
&\qquad+LR\cos(45^\circ+(m-1)\varphi)=T'\\
&mR^m\sin45^\circ+(m-1)AR^{m-1}\sin(45^\circ+\varphi)\\
&\qquad+(m-2)BR^{m-2}\sin(45^\circ+2\varphi)\\
&\quad+(m-3)CR^{m-3}\sin(45^\circ+3\varphi)+\text{usw.}\\
&\qquad+LR\sin(45^\circ+(m-1)\varphi)=U'
\end{aligned}
\]
Es ist klar, dass

I. $T$ aus den Teilen
\SourceLeaf{068}
\[
\begin{aligned}
&\frac{R^{m-1}}{m\BareRoot{2}}[R+mA\BareRoot{2}\, .\,\cos(45^\circ+\varphi)]\\
&+\frac{R^{m-2}}{m\BareRoot{2}}[RR+mB\BareRoot{2}\, .\,\cos(45^\circ+2\varphi)]\\
&+\frac{R^{m-3}}{m\BareRoot{2}}[R^3+mC\BareRoot{2}\, .\,\cos(45^\circ+3\varphi)]\\
&+\frac{R^{m-4}}{m\BareRoot{2}}[R^4+mD\BareRoot{2}\, .\,\cos(45^\circ+4\varphi)]\\
&+\text{usw.}
\end{aligned}
\]
besteht, von denen jeder bei jedem bestimmten reellen Wert von $\varphi$ positiv wird, wie man leicht erkennt. Daher erhält $T$ notwendig einen positiven Wert. Ebenso beweist man, dass auch $U,T',U'$ positiv werden, woraus folgt, dass auch $TT'+UU'$ notwendig eine positive Größe ist.

II. Für $r=R$ gehen die Funktionen $t,u,t',u'$ der Reihe nach in
\[
\begin{gathered}
T\cos(45^\circ+m\varphi)+U\sin(45^\circ+m\varphi)\\
T\sin(45^\circ+m\varphi)-U\cos(45^\circ+m\varphi)\\
T'\cos(45^\circ+m\varphi)+U'\sin(45^\circ+m\varphi)\\
T'\sin(45^\circ+m\varphi)-U'\cos(45^\circ+m\varphi)
\end{gathered}
\]
über, wie man durch Ausmultiplizieren leicht beweist. Daraus ergibt sich für die Funktion $tt'+uu'$ bei $r=R$ der Wert $=TT'+UU'$; dieser ist also eine positive Größe. Q. E. D.

Aus denselben Formeln folgt außerdem, dass der Wert der Funktion $tt+uu$ für $r=R$ gleich $TT+UU$, also positiv ist. Wir schließen daraus, dass für keinen Wert von $r$, der jede der Größen $mA\BareRoot{2},\BareRoot{(mB\BareRoot{2})},\BareRoot[3]{(mC\BareRoot{2})}$ usw. übertrifft, gleichzeitig $t=0,u=0$ gelten kann.

\ThirdArticle{2}
\textsc{Satz.} \textit{Innerhalb der Grenzen $r=0$ und $r=R$ sowie $\varphi=0$ und $\varphi=360^\circ$ gibt es gewiss solche Werte der Unbestimmten $r,\varphi$, für die gleichzeitig $t=0$ und $u=0$ gilt.}

\textit{Beweis.} Wir nehmen an, der Satz sei nicht wahr. Offenbar muss dann der Wert von $tt+uu$ für sämtliche Werte der Unbestimmten innerhalb der angegebenen Grenzen eine positive Größe sein, und daher muss der Wert von $y$ stets endlich bleiben. Wir betrachten das Doppelintegral
\[
\iint y\,\mathrm dr\,\mathrm d\varphi
\]
von $r=0$ bis $r=R$ und von $\varphi=0$ bis $\varphi=360^\circ$, das somit einen endlichen, vollständig bestimmten Wert erhält. Dieser Wert, den
\SourceLeaf{069}
wir mit $\Omega$ bezeichnen werden, muss derselbe sein, ob man zuerst nach $\varphi$ und dann nach $r$ oder in umgekehrter Reihenfolge integriert. Behandeln wir $r$ als konstant, so erhalten wir aber als \textit{unbestimmtes Integral}
\[
\int y\,\mathrm d\varphi=\frac{tu'-ut'}{r(tt+uu)}
\]
wie man durch Differentiation nach $\varphi$ leicht bestätigt. Eine Konstante ist nicht hinzuzufügen, wenn wir das Integral bei $\varphi=0$ beginnen lassen, da für $\varphi=0$ der Wert $\frac{tu'-ut'}{r(tt+uu)}=0$ ist. Da nun offenbar $\frac{tu'-ut'}{r(tt+uu)}$ auch für $\varphi=360^\circ$ verschwindet, ist das Integral $\int y\,\mathrm d\varphi$ von $\varphi=0$ bis $\varphi=360^\circ$ $=0$, während $r$ unbestimmt bleibt. Daraus folgt $\Omega=0$.

\Needspace{23\baselineskip}
Ebenso erhalten wir als unbestimmtes Integral, wenn wir $\varphi$ als konstant betrachten,
\[
\int y\,\mathrm dr=\frac{tt'+uu'}{tt+uu}
\]
wie man ebenso leicht durch Differentiation nach $r$ bestätigt. Auch hier ist keine Konstante hinzuzufügen, wenn das Integral bei $r=0$ beginnt. Nach dem im vorhergehenden Artikel Bewiesenen ist daher das Integral von $r=0$ bis $r=R$ $=\frac{TT'+UU'}{TT+UU}$ und somit nach dem dortigen Satz für jeden reellen Wert von $\varphi$ stets eine positive Größe. Daher ist auch $\Omega$, also der Wert des Integrals
\[
\int\frac{TT'+UU'}{TT+UU}\,\mathrm d\varphi
\]
von $\varphi=0$ bis $\varphi=360^\circ$, notwendig eine positive Größe\HistoricalNote{13}. Das ist absurd, da wir dieselbe Größe zuvor als $=0$ gefunden haben. Die Annahme kann also nicht bestehen, und damit ist die Wahrheit des Satzes erwiesen.

\ThirdArticle{3}
Die Funktion $X$ geht durch die Substitution $x=r(\cos\varphi+\sin\varphi\, .\,\BareRoot{-1})$ in $t+u\BareRoot{-1}$ über und durch die Substitution $x=r(\cos\varphi-\sin\varphi\, .\,\BareRoot{-1})$ in $t-u\BareRoot{-1}$. Wenn also für bestimmte Werte von $r,\varphi$, etwa für $r=g,\varphi=G$, gleichzeitig $t=0,u=0$ gilt (dass solche Werte existieren, wurde im vorhergehenden Artikel bewiesen), erhält $X$ durch jede der beiden Substitutionen
\SourceLeaf{070}
\[
x=g(\cos G+\sin G\, .\,\BareRoot{-1}),\quad x=g(\cos G-\sin G\, .\,\BareRoot{-1})
\]
den Wert 0 und ist daher identisch durch
\[
x-g(\cos G+\sin G\, .\,\BareRoot{-1}),\quad\text{sowie durch}\quad x-g(\cos G-\sin G\, .\,\BareRoot{-1})
\]
teilbar. Wenn weder $\sin G=0$ noch $g=0$ gilt, sind diese Teiler verschieden, und daher ist $X$ auch durch ihr Produkt
\[
xx-2g\cos G\, .\,x+gg
\]
teilbar. Wenn dagegen entweder $\sin G=0$, also $\cos G=\pm1$, oder $g=0$ gilt, sind jene Faktoren identisch, nämlich $=x\mp g$. Somit enthält die Funktion $X$ gewiss einen reellen Teiler zweiten oder ersten Grades. Da derselbe Schluss wiederum für den Quotienten gilt, lässt sich $X$ vollständig in solche Faktoren zerlegen. Q.~E.~D.

\ThirdArticle{4}
Obwohl wir die gestellte Aufgabe im Vorhergehenden bereits vollständig gelöst haben, wird es nicht überflüssig sein, noch einiges über die Schlussweise des Artikels 2 hinzuzufügen. Aus der Annahme, dass $t$ und $u$ für keine Werte der Unbestimmten $r,\varphi$ innerhalb der dort angegebenen Grenzen gleichzeitig verschwinden, sind wir zu einem unvermeidlichen Widerspruch gelangt und haben daraus die Unrichtigkeit der Annahme geschlossen. Dieser Widerspruch muss also wegfallen, wenn es tatsächlich Werte von $r,\varphi$ gibt, für die $t$ und $u$ gleichzeitig $=0$ werden. Zur näheren Erläuterung bemerken wir, dass für solche Werte $tt+uu=0$ gilt und somit $y$ selbst unendlich wird. Daher darf man das Doppelintegral $\iint y\,\mathrm dr\,\mathrm d\varphi$ nicht mehr als eine Größe mit bestimmtem Wert behandeln. Bezeichnen allgemein $\xi,\eta,\zeta$ unbestimmte Koordinaten von Punkten im Raum, so stellt das Integral $\iint y\,\mathrm dr\,\mathrm d\varphi$ das Volumen des Körpers dar, der zwischen den fünf Ebenen mit den Gleichungen
\[
\xi=0,\quad\eta=0,\quad\zeta=0,\quad\xi=R,\quad\eta=360^\circ
\]
und der Fläche mit der Gleichung $\zeta=y$ enthalten ist, wobei diejenigen Teile als negativ betrachtet werden, in denen die Koordinaten $\zeta$ negativ sind. Dabei wird jedoch stillschweigend vorausgesetzt, dass die sechste Fläche \textit{stetig} ist. Fällt diese Bedingung weg, indem $y$ unendlich wird, so kann jener Begriff durchaus seinen Sinn verlieren. In einem solchen Fall kann von einer Summierung des Integrals $\iint y\,\mathrm dr\,\mathrm d\varphi$ keine Rede sein, und es ist daher nicht verwunderlich, dass analytische Operationen, die man durch blindes Rechnen auf leere Begriffe anwendet, zu Widersinn führen.

\SourceLeaf{071}
Die Integration $\int y\,\mathrm d\varphi=\frac{tu'-ut'}{r(tt+uu)}$ ist nur insofern eine wirkliche Integration, also eine Summierung, als $y$ überall zwischen den Grenzen, über die sie sich erstreckt, eine endliche Größe ist. Sie ist dagegen widersinnig, wenn $y$ irgendwo zwischen diesen Grenzen unendlich wird. Wenn wir ein Integral wie $\int\eta\,\mathrm d\xi$, das allgemein die Fläche zwischen der Abszissenachse und der Kurve darstellt, deren Ordinate für die Abszisse $\xi$ $=\eta$ ist, nach den üblichen Regeln berechnen, ohne die Stetigkeit zu beachten, verwickeln wir uns sehr oft in Widersprüche. Setzen wir beispielsweise $\eta=\frac1{\xi\xi}$, so liefert die Analysis das Integral $=C-\frac1\xi$, das die Fläche richtig bestimmt, solange die Kurve stetig bleibt. Da diese Stetigkeit bei $\xi=0$ unterbrochen wird, erhält jemand, der unpassenderweise nach der Größe der Fläche von einer negativen bis zu einer positiven Abszisse fragt, von der Formel die widersinnige Antwort, sie sei negativ. Was diese und ähnliche Paradoxien der Analysis bedeuten, werden wir bei anderer Gelegenheit ausführlicher verfolgen.

\Needspace{22\baselineskip}
Hier sei nur noch eine Bemerkung hinzugefügt. Wenn Fragen, die in gewissen Fällen widersinnig werden können, \textit{ohne Einschränkung} gestellt werden, hilft sich die Analysis sehr oft dadurch, dass sie die Antwort zum Teil unbestimmt lässt. Für den Wert des Integrals $\iint y\,\mathrm dr\,\mathrm d\varphi$ von $r=e$ bis $r=f$ und von $\varphi=E$ bis $\varphi=F$ erhält man nämlich, wenn der Wert von $\frac ut$
\[
\begin{aligned}
&\text{für }r=e,\ \varphi=E\quad\text{mit}\quad\theta\\
&\qquad r=e,\ \varphi=F\quad .\quad .\quad .\quad .\quad .\quad .\quad .\quad\theta'\\
&\qquad r=f,\ \varphi=E\quad .\quad .\quad .\quad .\quad .\quad .\quad .\quad\theta''\\
&\qquad r=f,\ \varphi=F\quad .\quad .\quad .\quad .\quad .\quad .\quad .\quad\theta'''
\end{aligned}
\]
bezeichnet wird, durch analytische Operationen leicht
\[
\mathrm{Arc.\ tang}\,\theta-\mathrm{Arc.\ tang}\,\theta'-\mathrm{Arc.\ tang}\,\theta''+\mathrm{Arc.\ tang}\,\theta'''
\]
Tatsächlich kann das Integral nur dann einen bestimmten Wert haben, wenn $y$ innerhalb der angegebenen Grenzen stets endlich bleibt. Dieser Wert ist zwar in der angegebenen Formel enthalten, wird durch sie aber noch nicht vollständig bestimmt, da Arc. tang. eine mehrwertige Funktion ist. Man wird gesondert durch andere Überlegungen (die allerdings nicht schwierig sind) entscheiden müssen, welche Werte der Funktion im bestimmten Fall gerade zu verwenden sind. Wird dagegen $y$ irgendwo innerhalb der angegebenen Grenzen unendlich, so ist die Frage nach dem Wert des Integrals $\iint y\,\mathrm dr\,\mathrm d\varphi$ widersinnig. Besteht man trotzdem darauf, der Analysis eine Antwort abzuzwingen, so wird sie je nach der Verschiedenheit der Methoden bald diese, bald jene geben; jede einzelne dieser Antworten wird jedoch in der zuvor angegebenen allgemeinen Formel enthalten sein.
\begin{center}\rule{3cm}{0.3pt}\end{center}
